\section*{Solutions: Exam 2 - Abstract Algebra}

\subsection{Problem 1a: Proving $(\mathbb{Z}, *)$ is a Group}

\begin{solution}
Given: $a * b = a + b + 2026$ for all $a, b \in \mathbb{Z}$

We need to verify:
\begin{enumerate}
    \item \textbf{Closure:} For all $a, b \in \mathbb{Z}$, $a * b = a + b + 2026 \in \mathbb{Z}$ ✓
    
    \item \textbf{Associativity:} 
    \begin{align*}
        (a * b) * c &= (a + b + 2026) * c\\
        &= (a + b + 2026) + c + 2026\\
        &= a + b + c + 2 \cdot 2026
    \end{align*}
    
    \begin{align*}
        a * (b * c) &= a * (b + c + 2026)\\
        &= a + (b + c + 2026) + 2026\\
        &= a + b + c + 2 \cdot 2026
    \end{align*}
    
    So $(a * b) * c = a * (b * c)$ ✓
    
    \item \textbf{Identity element:} We need $e$ such that $a * e = a$ for all $a \in \mathbb{Z}$
    
    $a * e = a + e + 2026 = a$
    
    Thus: $e = -2026$
    
    Check: $a * (-2026) = a + (-2026) + 2026 = a$ ✓
    
    \item \textbf{Inverse element:} For each $a \in \mathbb{Z}$, we need $a' \in \mathbb{Z}$ such that $a * a' = e = -2026$
    
    $a + a' + 2026 = -2026$
    
    Thus: $a' = -4052 - a \in \mathbb{Z}$
    
    Check: $a * (-4052 - a) = a + (-4052 - a) + 2026 = -2026$ ✓
\end{enumerate}

\textbf{Conclusion:} $(\mathbb{Z}, *)$ satisfies all group axioms, so it is a group.

\end{solution}

\subsection{Problem 1b: Is $X = \{4k+2 | k \in \mathbb{Z}\}$ a Subgroup?}

\begin{solution}
$X$ is the set of all even integers of the form $4k+2$ (i.e., 2, 6, 10, -2, -6, ...).

For $X$ to be a subgroup, we need:
\begin{enumerate}
    \item $X \neq \emptyset$ (clearly true)
    \item For $a, b \in X$, the inverse $a' \in X$ and $a * a' \in X$ (subgroup criterion)
\end{enumerate}

\textbf{Check identity element:} The identity is $e = -2026$

Is $-2026 \in X$? We need $-2026 = 4k + 2$ for some $k \in \mathbb{Z}$

$-2026 = 4k + 2 \Rightarrow 4k = -2028 \Rightarrow k = -507$

Yes! $-2026 = 4(-507) + 2 \in X$ ✓

\textbf{Check closure:} Take $a = 2, b = 6 \in X$

$a * b = 2 + 6 + 2026 = 2034 = 4(508) + 2 \in X$ ✓

\textbf{Check inverses:} For $a = 2$, the inverse is $a' = -4052 - 2 = -4054 = 4(-1014) + 2 \in X$ ✓

Actually, let me verify more generally: If $a = 4k + 2$, then
$$a' = -4052 - a = -4052 - (4k + 2) = -4054 - 4k = 4(-1013.5 - k)$$

Wait, $-4054 = 4(-1013.5)$, which is not an integer coefficient. Let me recalculate:

$-4054 / 4 = -1013.5$, so $a' = -4054 - 4k$ doesn't have the form $4m + 2$.

Let me verify: $-4054 = 4(-1014) + 2$? 
$4(-1014) + 2 = -4056 + 2 = -4054$ ✓

So $a' = 4(-1014) - 4k = 4(-1014 - k) + 2 \in X$ ✓

\textbf{Conclusion:} $X$ is a subgroup of $(\mathbb{Z}, *)$

\end{solution}

\subsection{Problem 2: Group Homomorphism}

\begin{solution}
Given: $\varphi: X \to Y; x^l \mapsto (y^k)^l$ where $X$ is cyclic order $m$, $Y$ is cyclic order $n$

Hypothesis: $km \equiv 0 \pmod{n}$ (i.e., $n | km$)

We need to prove $\varphi$ is a homomorphism: $\varphi(x^a \cdot x^b) = \varphi(x^a) \cdot \varphi(x^b)$

\textbf{Verification:}

In $X$: $x^a \cdot x^b = x^{a+b}$ (multiplicatively)

$$\varphi(x^{a+b}) = (y^k)^{a+b} = y^{k(a+b)}$$

$$\varphi(x^a) \cdot \varphi(x^b) = (y^k)^a \cdot (y^k)^b = y^{ka} \cdot y^{kb} = y^{k(a+b)}$$

So $\varphi(x^{a+b}) = \varphi(x^a) \cdot \varphi(x^b)$ ✓

\textbf{Well-definedness:} We need to ensure the mapping is well-defined in $Y$.

For $\varphi$ to be well-defined, if $x^l = x^{l'}$ in $X$, then $(y^k)^l = (y^k)^{l'}$ in $Y$.

$x^l = x^{l'}$ in $X$ means $l \equiv l' \pmod{m}$

$(y^k)^l = (y^k)^{l'}$ in $Y$ means $kl \equiv kl' \pmod{n}$

Given $n | km$, we have $kl \equiv kl' \pmod{n}$ whenever $l \equiv l' \pmod{m}$ ✓

\textbf{Conclusion:} $\varphi$ is a group homomorphism.

\end{solution}

\subsection{Problem 3a: $\mathbb{Z}(\sqrt{2})$ is a Ring}

\begin{solution}
$\mathbb{Z}(\sqrt{2}) = \{a + b\sqrt{2} | a, b \in \mathbb{Z}\}$ with standard real addition and multiplication.

\textbf{Verification of ring axioms:}

\begin{enumerate}
    \item \textbf{Closure under addition:} $(a_1 + b_1\sqrt{2}) + (a_2 + b_2\sqrt{2}) = (a_1+a_2) + (b_1+b_2)\sqrt{2} \in \mathbb{Z}(\sqrt{2})$ ✓
    
    \item \textbf{Additive identity:} $0 = 0 + 0\sqrt{2} \in \mathbb{Z}(\sqrt{2})$ ✓
    
    \item \textbf{Additive inverses:} $-(a + b\sqrt{2}) = (-a) + (-b)\sqrt{2} \in \mathbb{Z}(\sqrt{2})$ ✓
    
    \item \textbf{Associativity of addition:} Inherited from $\mathbb{R}$ ✓
    
    \item \textbf{Commutativity of addition:} Inherited from $\mathbb{R}$ ✓
    
    \item \textbf{Closure under multiplication:} 
    $(a_1 + b_1\sqrt{2})(a_2 + b_2\sqrt{2}) = a_1a_2 + 2b_1b_2 + (a_1b_2 + a_2b_1)\sqrt{2} \in \mathbb{Z}(\sqrt{2})$ ✓
    
    \item \textbf{Multiplicative identity:} $1 = 1 + 0\sqrt{2} \in \mathbb{Z}(\sqrt{2})$ ✓
    
    \item \textbf{Associativity of multiplication:} Inherited from $\mathbb{R}$ ✓
    
    \item \textbf{Distributivity:} Inherited from $\mathbb{R}$ ✓
    
    \item \textbf{Commutativity of multiplication:} Inherited from $\mathbb{R}$ ✓
\end{enumerate}

\textbf{Conclusion:} $\mathbb{Z}(\sqrt{2})$ is a commutative ring with unity.

\end{solution}

\subsection{Problem 3b: Is $\mathbb{Z}(\sqrt{2})$ a Field?}

\begin{solution}
For $\mathbb{Z}(\sqrt{2})$ to be a field, every nonzero element must have a multiplicative inverse in $\mathbb{Z}(\sqrt{2})$.

\textbf{Counterexample:} Consider $a = 2 \in \mathbb{Z}(\sqrt{2})$

The multiplicative inverse of 2 would be $2^{-1} = \frac{1}{2}$

Is $\frac{1}{2} \in \mathbb{Z}(\sqrt{2})$? We need integers $p, q$ such that:
$$\frac{1}{2} = p + q\sqrt{2}$$

This implies $1 = 2p + 2q\sqrt{2}$

Comparing rational and irrational parts:
- Rational: $2p = 1 \Rightarrow p = 1/2 \notin \mathbb{Z}$
- Irrational: $2q = 0 \Rightarrow q = 0$

This is a contradiction, so $2^{-1} \notin \mathbb{Z}(\sqrt{2})$.

\textbf{Conclusion:} $\mathbb{Z}(\sqrt{2})$ is NOT a field. However, it is an integral domain (as it's a subring of $\mathbb{R}$).

\end{solution}

\subsection{Problem 4: Field, Integral Domain, and Properties}

\begin{solution}

\textbf{a) Definitions:}

\textbf{Field:} A commutative ring $F$ with unity $1 \neq 0$ where every nonzero element has a multiplicative inverse.

\textbf{Integral Domain:} A commutative ring $D$ with unity $1 \neq 0$ with no zero divisors (i.e., $ab = 0 \Rightarrow a = 0$ or $b = 0$).

\textbf{b) Every field is an integral domain:}

Let $F$ be a field. Suppose $ab = 0$ for $a, b \in F$.

If $a \neq 0$, then $a^{-1}$ exists. Thus:
$$b = a^{-1}(ab) = a^{-1} \cdot 0 = 0$$

So no zero divisors exist. Hence $F$ is an integral domain. ✓

\textbf{c) Is $\mathbb{Z}^n$ ($n > 1$) an integral domain?}

\textbf{No.} $\mathbb{Z}^n$ has zero divisors.

\textbf{Counterexample:} In $\mathbb{Z}^2$, let $a = (1, 0)$ and $b = (0, 1)$.

Both $a, b \neq (0, 0)$, but $ab = (1 \cdot 0, 0 \cdot 1) = (0, 0)$.

Hence $\mathbb{Z}^2$ is not an integral domain.

\end{solution}

\subsection{Problem 5a: Eisenstein's Criterion}

\begin{solution}
Polynomial: $f(x) = x^3 - 3x + 1 \in \mathbb{Q}[x]$

\textbf{Eisenstein's Criterion:} If for prime $p$:
\begin{enumerate}
    \item $p | a_i$ for all $i < n$ (all non-leading coefficients)
    \item $p \nmid a_n$ (leading coefficient)
    \item $p^2 \nmid a_0$ (constant term)
\end{enumerate}
Then $f(x)$ is irreducible over $\mathbb{Q}$.

Try $p = 3$:
- Coefficients: $1, 0, -3, 1$
- $3 | 0, 3 | (-3)$ (non-leading coefficients) ✓
- $3 \nmid 1$ (leading coefficient) ✓
- $3^2 = 9 \nmid 1$ (constant term) ✓

\textbf{Conclusion:} $f(x) = x^3 - 3x + 1$ is irreducible in $\mathbb{Q}[x]$ by Eisenstein's criterion with $p = 3$.

\end{solution}

\subsection{Problem 5b: Polynomial Division in $\mathbb{Z}_7[x]$}

\begin{solution}
$f = x^3 + \overline{p}x + \overline{5}$ and $g = x^2 + \overline{5}x + \overline{6}$ in $\mathbb{Z}_7[x]$

For $g | f$, we need $f = gq + r$ where $\deg(r) < 2$.

Divide: $f = g(x - 5) + (ax + b)$

$$x^3 + \overline{p}x + \overline{5} = (x^2 + \overline{5}x + \overline{6})(x - \overline{5}) + ax + b$$

Expand RHS:
$$(x^2 + \overline{5}x + \overline{6})(x - \overline{5}) = x^3 - \overline{5}x^2 + \overline{5}x^2 - \overline{25}x + \overline{6}x - \overline{30}$$

$$= x^3 + (\overline{6} - \overline{25})x - \overline{30}$$

In $\mathbb{Z}_7$: $-25 \equiv -25 + 28 = 3 \pmod{7}$ and $-30 \equiv -30 + 35 = 5 \pmod{7}$

$$= x^3 + (\overline{6} + \overline{3})x + \overline{5} = x^3 + \overline{2}x + \overline{5}$$

For $g | f$, we need remainder = 0:
$$x^3 + \overline{p}x + \overline{5} = x^3 + \overline{2}x + \overline{5}$$

Thus: $\overline{p} = \overline{2}$, so $\boxed{p \equiv 2 \pmod{7}}$

\end{solution}

\newpage
